%!TeX program = xelatex
\documentclass{SYSUReport}
\usepackage{tabularx} % 在导言区添加此行
\usepackage{float}
% 根据个人情况修改
\headl{}
\headc{}
\headr{并行程序设计与算法实验}
\lessonTitle{并行程序设计与算法实验}
\reportTitle{Lab8-并行多源最短路径搜索}
\stuname{xxx}
\stuid{xxx}
\inst{计算机学院}
\major{xxx}
\date{2025年5月14日}

\begin{document}

% =============================================
% Part 1: 封面
% =============================================
\cover
\thispagestyle{empty} % 首页不显示页码
\clearpage

% % =============================================
% % Part 4： 正文内容
% % =============================================
% % 重置页码，并使用阿拉伯数字
% \pagenumbering{arabic}
% \setcounter{page}{1}

%%可选择这里也放一个标题
%\begin{center}
%    \title{ \Huge \textbf{{标题}}}
%\end{center}

\section{实验目的}
\begin{itemize}
    \item 评估并行化最短路径算法的性能。
    \item 探索不同最短路径算法的并行适应性。
\end{itemize}
\section{实验内容}
实现并行的多源最短路径搜索，具体要求如下：
\begin{itemize}
    \item 使用 OpenMP、Pthreads 或 MPI 中的任意一种并行编程模型来实现。

 \textbf{你选择的并行框架是：}xxx
    \item 选用任意一种最短路径算法作为基础进行并行化，如Bellman-Ford、Dijkstra（针对每个源点执行）、Floyd-Warshall 或 Johnson 算法。
    
\textbf{你选用的算法是：}xxx
\item 简要描述所选算法的实现要点。

\textbf{回答：}xxx
\end{itemize}
\section{实验结果与分析}
\subsection{实验配置简述}
\begin{itemize}
    \item \textbf{所用算法：} [请在此处填写算法名称]
    \item \textbf{并行框架：} [请在此处填写 OpenMP/Pthreads/MPI]
    \item \textbf{并行方式：} [请在此处简述核心并行策略，例如：外层循环并行化、数据划分方式等]
    \item \textbf{测试数据特征：}
        \begin{itemize}
            \item 数据集1：节点数量 [V1], 平均度数 [D1] (或边数量 [E1])
            \item 数据集2 (若有)：节点数量 [V2], 平均度数 [D2] (或边数量 [E2])
            \item (可添加更多数据集)
        \end{itemize}
\end{itemize}
\textbf{数据集1 (节点数量: [V1], 平均度数: [D1]) 的性能数据：}
\begin{table}[H]
  \centering
  \caption{数据集1 ([V1]节点, 平均度数[D1]) 上 [算法名称] 的并行性能 (串行时间 $T_s$: [填写T1时间] 秒)}
  \label{tab:results_dataset1}
  \begin{tabular}{cccc}
    \toprule
    \textbf{线程数量 ($p$)} & \textbf{运行时间 $T_p$ (秒)} & \textbf{加速比 $S_p = T_s / T_p$} & \textbf{并行效率 $E_p = S_p / p$} \\
    \midrule
    1   & [时间] & 1.00     & 1.00      \\
    2   & [时间] & [计算值] & [计算值]  \\
    4   & [时间] & [计算值] & [计算值]  \\
    8   & [时间] & [计算值] & [计算值]  \\
    16  & [时间] & [计算值] & [计算值]  \\
    % 可根据实际测试的线程数调整或增加行
    \bottomrule
  \end{tabular}
\end{table}

\textit{(如果测试了多个数据集，请为每个数据集复制并填写类似的表格)}

% \textbf{针对 数据集2 (节点数量: [V2], 平均度数: [D2]) 的性能数据：}
% \begin{table}[htbp]
%   \centering
%   \caption{数据集2 ([V2]节点, 平均度数[D2]) 上 [算法名称] 的并行性能 (串行时间 $T_s$: [填写T1时间] 秒)}
%   \label{tab:results_dataset2}
%   \begin{tabular}{cccc}
%     \toprule
%     \textbf{线程数量 ($p$)} & \textbf{运行时间 $T_p$ (秒)} & \textbf{加速比 $S_p = T_s / T_p$} & \textbf{并行效率 $E_p = S_p / p$} \\
%     \midrule
%     1   & [时间] & 1.00     & 1.00      \\
%     2   & [时间] & [计算值] & [计算值]  \\
%     4   & [时间] & [计算值] & [计算值]  \\
%     8   & [时间] & [计算值] & [计算值]  \\
%     16  & [时间] & [计算值] & [计算值]  \\
%     \bottomrule
%   \end{tabular}
% \end{table}
\subsection{并行性能分析}
根据你的实验数据，结合你选择的算法、并行框架、并行方式以及测试数据的特征（节点数量、平均度数等），分析程序的并行性能。可辅以图表（如加速比曲线、效率曲线）进行更直观的分析。
\begin{itemize}
    \item 运行时间随线程数增加的变化趋势。

    \textbf{回答：}
    \item 加速比的变化情况。

        \textbf{回答：}

    \item 并行效率的变化情况，分析影响并行效率的主要因素。

        \textbf{回答：}

    \item 不同数据特征（如节点数量、平均度数）对并行性能的影响。

        \textbf{回答：}

\end{itemize}

注：实验报告格式参考本模板，可在此基础上进行修改；实验代码以zip格式另提交；最终提交内容包括实验报告(pdf格式)和实验代码(zip压缩包格式)
\end{document}